\noindent\textbf{Por qué se usa este procedimiento.} Calcular primero $G=E/[2(1+\nu)]$ para el material isótropo. La ley de cortante $\tau=G\gamma$ se despeja como $\gamma=\tau/G$. El ángulo de distorsión sale adimensional, aproximadamente en radianes.

\begin{enumerate}

\item El módulo transversal es $G = E/[2(1+\nu)]$.

\item La distorsión es $\gamma = \tau/G$.

\end{enumerate}

\noindent\textbf{Operaciones y resultados intermedios.}

\begin{enumerate}

\item Módulo transversal.

\[\fitmath{y_{1}=\frac{71700}{2\cdot \left(1+0{,}33\right)}\approx 26954{,}89\,\mathrm{MPa}}\]

\item Distorsión.

\[\fitmath{y_{2}=\frac{50}{\frac{71700}{2\cdot \left(1+0{,}33\right)}}\approx 0{,}001854951\,\mathrm{rad}}\]

\end{enumerate}
